\section{Local contracts and pass correspondence}
\label{sec:pass-contracts}
Write $\mathsf{WF}_L(s)\Downarrow A$ for successful local checking with analysis
$A$ (unit for checks returning no analysis). Write
$\mathsf{check}_f(s,t,\chi)\Downarrow\ok$ for the separately callable pass
correspondence check. These are executable structural checks. They do not
establish Rust expression semantics, representation laws, or equality of
arbitrary program meanings. Rules below summarize the predicates defined here;
each implementation rule is registered at the function that checks them.

\subsection{Guard recognition and short-circuit refinement}
The CQ parser recognizes the trait operations in CQ-GuardRecognize after
collecting all query binders. Its output contains relational clauses and only
residual opaque guards. The CQ contract rejects a recognizable, unsplit guard.
Let $e_1,\ldots,e_n$ be the leaves of the top-level Rust conjunction, with
parentheses removed only for this analysis. Partition them into recognized
total predicates $A$ and opaque leaves $F$, preserving each written sequence.
Mixed disjunctions and negations are single opaque leaves.
\begin{mathpar}
\Rule{CQ-GuardSplit}{e=e_1\mathbin{\&\&}\cdots\mathbin{\&\&}e_n\\
 A=\text{recognized leaves}\quad F=\text{remaining leaves in written order}}
 {\mathsf{if}\ e\rightsquigarrow \mathsf{select}_{\bigwedge A};
   \mathsf{guard}_{\bigwedge F}\quad
   \sem{\mathsf{select}_{\bigwedge A};\mathsf{guard}_{\bigwedge F}}
   \sqsupseteq\sem{\mathsf{if}\ e}}
\end{mathpar}
The residual is omitted when $F$ is empty, and otherwise remains at the
written guard position. The recognized predicates may then move earlier
under the existing relational pass contracts.
\begin{proof}
Rust evaluates $f(x)$ in \code{a \&\& f(x)} exactly where $a$ holds. Selecting
by total relational $a$ earlier yields exactly that domain. In
\code{f(x) \&\& a}, Rust evaluates $f$ on a superset of that domain, so selecting
$a$ earlier only shrinks it. Iterate over the conjunction while retaining
opaque leaves in their written order. Removing a total relational leaf cannot
introduce a panic: on every input where the written evaluation returns, either
some relational leaf rejects the input before any residual evaluation, or all
relational leaves hold and the residual evaluates the same opaque leaves in
the same order. Returned sets agree. Earlier relational selection may avoid a
written panic, as permitted by $\bot+\mathrm{Set}$ refinement. Extracting
through a mixed disjunction or negation lacks this argument and is forbidden.
\end{proof}

\subsection{Language-local contracts}
\paragraph{CQ and CQ-NF.}
$\mathsf{CQScope}(s)$ consists of the CQ clause rules in
Section~\ref{sec:languages}, plus declaration and result checks: relation and
column names are unique and declared relations have positive arity;
declarations contain no attributes; declaration order fixes host row layout;
an optional parenthesized leaf expression is preserved as Rust syntax; body atoms name declared inputs at the declared arity;
lets cannot shadow bound logical names; the body contains an atom; the head has
positive arity. A declaration named by the head is treated as the result and
must have its arity and cannot bind a leaf value. Without that declaration, every head argument must be a
bound bare variable with a known, non-inference symbolic type. With a declared
result, CQ delegates head expression typing to Rust; NF still requires every
head variable to be body-bound after normalization. Blocks check common
interface types and reject escaped branch-local names.
Negation and aggregation parse but are rejected by CQ checking. An opaque
\code{name! \{tokens\}} clause also parses and round-trips, but CQ checking
rejects it with the name of the unexpanded extension. Visible logical
uses are checked in written scope; opaque syntax remains Rust's responsibility.
Central comparison terms admit only bound query names, literals, or captured
identifiers. Recognizer shape and known Eq/Ord type failures are checked locally;
trait laws and semantic admissibility remain trusted premises.
\begin{mathpar}

\Rule{WF-CQ}{\mathsf{CQScope}(s)}{\mathsf{WF}_{CQ}(s)\Downarrow\ok}

\Rule{WF-NF}{\mathsf{WF}_{CQ}(s)\quad\text{every nested atom has distinct bare arguments}\\ \text{head arguments are bound variables; known inadmissible types rejected}}
 {\mathsf{WF}_{NF}(s)\Downarrow\ok}
\end{mathpar}
NF is a restriction of the same \code{cq::Program}, checked by
\code{normal::check}, not another Rust AST type. Its head is an ordered tuple
of variables and may repeat them; computed heads have already become binds.

\paragraph{VO.}
$\mathsf{Schedule}(t,A)$ checks each subatom's relation, clause, arity,
original column names and complete correlated-row coverage. Dropped columns
must be unbound. Each relational identity and central constraint appears exactly once, uses
only bound query variables and binds nothing. Block steps recursively check
operator kind, paths, interface coverage and each arm's schedule. Rust lets and guards keep one identical
written chain; every free-variable binder and every source Before relational
clause must precede them. Relational clauses can move earlier. A semijoin reduction names a different atom or Rust-free derived block and
lists the exact shared equality classes in its support's column positions.
Ordinary targets are reduced before their first subatom; a target already read
while supplying a finite block domain instead restricts the retained prefix. No opaque
Rust guard is a reduction. VO does not require an acyclic reduction graph.
\begin{mathpar}
\Rule{D-Constraint}{\vars(C)\subseteq B}
 {\mathsf{requires}(C)=\vars(C),\quad\mathsf{provides}(C)=\emptyset}
\end{mathpar}
\[
\Rule{WF-VO}{\mathsf{WF}_{NF}(t.cq)\quad\mathsf{Schedule}(t,A)}
 {\mathsf{WF}_{VO}(t)\Downarrow A}.
\]

\paragraph{RA and IX.}
$\mathsf{RAHeading}(a,A)$ means the local operator rules in
Section~\ref{sec:languages}, valid declarations, fresh result names, earlier
operands, reachable definitions, visible expression-use checks, and a defined
output. An all-bare output has exactly its input heading; a computed output may
use a subset. Analysis carries bindings, symbolic types and computed-column
provenance, without Rust type unification. For IX let $D(a)$ contain the exact stream descriptors of joins/semijoins,
anti stream descriptors on its explicit free key, const descriptors of
point/membership conjunctions, and range descriptors of constant/parameter
order conjunctions on one column.
\begin{mathpar}

\Rule{WF-RA}{\mathsf{RAHeading}(a,A)}{\mathsf{WF}_{RA}(a)\Downarrow A}

\Rule{WF-IX}{\mathsf{WF}_{RA}(t.ra)\quad\text{each index ends in a semicolon}\\
 \text{unique index names and descriptors}\quad\operatorname{descs}(t)\subseteq D(t.ra)\\
 \text{each descriptor names a result with strictly increasing in-range ordinals}}
 {\mathsf{WF}_{IX}(t)\Downarrow\ok}
\end{mathpar}
IX local validity permits partial demand coverage and arbitrary index names.
First-use ordering, complete coverage and \code{iN} names belong to the pass
correspondence rule below.

\paragraph{Physical.}
The program is standalone. $\mathsf{PHeading}(t,A)$ checks declarations and
unique results, earlier operands, total one-to-one renames, fresh extension
columns, subset projections, visible expression uses, defined output and
reachability. Probe/merge/semi keys are duplicate-free lists whose set is exactly
the operands' heading intersection; semi keeps the left heading, probe/merge
take the union. Anti keeps the left heading and checks its explicit key is
a duplicate-free subset of the common heading. Union requires at least two
operands with equal headings and consistent known interface types; its block
origin checks free/interface, prefix/domain and witness metadata. Range checks a nonempty conjunction of constant/parameter order bounds on
exactly one available Ord column. Host layout and type choose direct access
or scan/filter later; no integer recognizer or Rust guard is involved. Shared algebra operators
have their local heading effects. Analysis derives rows, source identifiers,
symbolic types, computed columns, hosts/renames, keyed/merge/range reads, scans,
central boolean conditions, modes and canon flags. Output expressions
reject visible unavailable attributes; an all-bare head's dependency set must
exactly equal the returned heading, while a computed head may use a subset.
Let $\mathsf{reads}(r)$ count operand occurrences and the return use.
Here reusableHost means a noncanonical host rename; linearOrFused means
exactly one read or only shared prefix reads within a checked fused group.
\begin{mathpar}

\Rule{WF-Physical}{\mathsf{PHeading}(t,A)\quad\mathsf{PMode}(t,A)}
 {\mathsf{WF}_{Physical}(t)\Downarrow A}

\Rule{PH-Pipe}{\mathsf{mode}(r)=pipe\\
 \mathsf{reusableHost}(r)\ \lor\
 (\mathsf{linearOrFused}(r)\land r\notin\dom A.key\_reads)}
 {\mathsf{PModeOK}(r)}

\Rule{PH-Let}{\mathsf{mode}(r)\in\{let,exists\}\quad\mathsf{reads}(r)>0}
 {\mathsf{PModeOK}(r)}

\end{mathpar}
$\mathsf{PMode}$ requires the applicable rule for every result. An exists mode requires an empty heading and at least one read; it consumes
the full stream into retained existence state. A merge-side read can satisfy
a pipe's single-use rule; Structures discharges its ordering premise.
\code{probe smallest} requires an adjacent admissible intersection. Lazy
semi/anti requires a reusable pipe host with exclusively lazy membership
reads on that key. A witness union has an empty interface (or an unannotated
empty heading); a fused union requires its recorded prefix and every arm to
depend on it. Physical rejects cross-arm canon on statically disjoint unions.
Physical checks neither IX demand membership nor canonical return evidence.

\paragraph{Structures.}
$\mathsf{Entries}(t,A)$ validates its embedded Physical and records built
entries over retained results, per-reader host fallbacks, retained collections,
and canonical-boundary/computed-head dedup types. Built columns belong to the
heading and never repeat across levels; duplicate concatenated keys fail.
A built key answers the union of a prefix of its levels.
Write $\mathsf{hostRead}(r)$ when $r$ is a reusable pipe rename and
$\mathsf{built}(r,C)$ for a matching level prefix. Host entries record fallback types, not native trait claims.
\begin{mathpar}

\Rule{STR-Key}{\mathsf{hostRead}(r)\lor\mathsf{built}(r,C)}
 {\mathsf{answers}(r,C,probe/semi)}

\Rule{STR-Range}{\text{constant window on }r,c}
 {\mathsf{answers}(r,\{c\},range)\text{ by direct range or scan/filter; no build}}

\Rule{STR-Merge}{\mathsf{hostRead}(r)\lor\mathsf{built}(r,C)\\
 \text{or }(\mathsf{mode}(r)=pipe\land\mathsf{operator}(r)=merge[C](l,s))}
 {\mathsf{answers}(r,C,merge)}

\Rule{WF-Structures}{\mathsf{Entries}(t,A)\quad\text{all Physical reads answered}\\
 \mathsf{unique}(t.returned)\text{ inferred or explicit return canon}}
 {\mathsf{WF}_{Structures}(t)\Downarrow A}
\end{mathpar}
Uniqueness starts Unit at the empty key; host renames establish no key.
Join requires both input keys and unions them; semi/anti keeps the left key;
union establishes no inferred key; select/point/range/guard/bind preserves it; project preserves it only when no key column
is dropped. A canon flag establishes the full heading through its recorded
dedup type. A built structure alone proves no uniqueness. The return must
have a key or \code{return canon}. Every definition canon needs a dedup
entry; return canon needs a head dedup entry, which deduplicates answer values.
Computed heads also need that entry when they can collapse distinct
valuations. Retained scans, multiple unfiltered builds and exists state need
recorded row storage; host renames need scan and keyed-reader fallbacks.
Rust checks the traits of built structures and direct host reads;
trait laws remain representation assumptions, not syntactic conclusions.

Filtered entries record \code{payload[columns] constants{column = term}}.
Only point/range selections or unions can use them. For a selection, a single
shared one-level build has no row scan or retained-row entry. Noncorrelated
unions may have multiple direct indices. Every build records the same
projection; readers must be stream-supplied \code{probe first}, semi or anti
on a recorded level prefix. Payload plus exactly the singleton constants proved
by the selection must cover downstream keys, predicates, Rust free variables
and answers. Constant columns are reconstructed, not stored. These checks
derive emitted layouts; an opaque Rust payload conservatively retains its
whole visible environment and a bind is evaluated even if its outputs are unused.

\subsection{Same-language rewrite checks}
CQ rewrites are AST identity. VO requires unchanged CQ and exact equality
of non-reduction step erasures recursively, and preserves each existing
reduction while allowing added reductions subject to the target coverage
and Before contract. It does not permit rescheduling non-reduction steps. RA prime rewrites
use the same fixed-occurrence structural checker as elaboration, with unchanged
declarations and head. IX requires unchanged RA and locally valid demands;
it does not itself require complete first-use demands. Physical rewrites
preserve named logical operators, operands, declarations, head and returned
result, admitting changes to execution choices and modes and earlier
relational scheduling while retaining branch-scoped Before facts. This
standalone check has no IX source and does not establish final canonicality.
Structures rewrites check both local contracts and the Physical rewrite
relation; concrete representations can change when read coverage still holds.
Rust-expression rewrites in the generic protocol are AST identity; raw-token
surface rewrites have no semantic contract. Language-changing correspondence
checks below are separate from these rewrite contracts.

\subsection{Pass correspondence checks}
Each complete pass check requires successful local checks of its source
and target; equalities are the AST equalities used by the code, unless a rule
explicitly compares normalized names or descriptors. These checks do not call
\code{compile} to compare its output.

\paragraph{Ordering.}
Define $\mathsf{OrderObligations}(s,t,\chi)$ as: unchanged CQ; no reductions;
at least one subatom per atom clause; Coarse has
exactly one full subatom per atom; Fine introduces at most one new variable per
subatom. Under Written, introducing subatom indices are nondecreasing (all
Coarse subatoms count; in Fine only those introducing a name count).
Written preserves the complete NF body order (splitting atoms in place);
optimized guard/let placement is constrained by VO.
Prefix and Elimination rankings are construction choices checked by fixtures,
not correspondence obligations.
\[
\Rule{CK-Order}{\mathsf{WF}_{NF}(s)\quad\mathsf{WF}_{VO}(t)\\
 \mathsf{OrderObligations}(s,t,\chi)}
 {\mathsf{check}_{order}(s,t,\chi)\Downarrow\ok}.
\]

\paragraph{Reduction.}
$\mathsf{ReduceObligations}$ requires a reduction-free source, unchanged CQ,
and exact equality of surviving steps after erasing semijoin reductions.
Constraints and opaque Rust steps remain present. The legacy guards option
changes no Rust step. Under semijoins=None, reductions are forbidden; under
Tree, the source is acyclic, every directed join-tree edge occurs once, all
upward edges precede downward edges, and each edge respects parent/child
dependencies. Hypergraph additionally requires each nonempty atom-to-block
cut and uses Tree obligations only when acyclic; branch checks recurse. Sibling order is free. Local VO checks shared columns and
placement (including the finite-domain retained-prefix exception).
\[
\Rule{CK-Reduce}{\mathsf{WF}_{VO}(s)\quad\mathsf{WF}_{VO}(t)\\
 \mathsf{ReduceObligations}(s,t,\chi)}
 {\mathsf{check}_{reduce}(s,t,\chi)\Downarrow\ok}.
\]

\paragraph{Elaboration.}
The independent checker builds a schedule from the supplied VO, without
running elaboration or discovering occurrence mappings. It compares fixed
atom paths and positional column bindings, equality closure, predicate shape
with fixed column arguments, named reduction support, Rust payloads/bindings,
output heading and head. DAG traversal assigns facts to their first observing
Rust region; each source Before fact must occur no later in the target.
Independent projected views of one occurrence require exact interned structural
wiring, retaining witness correlation. This intentionally rejects some
semantically equal plans. Annotated blocks are checked through symbolic prefix
views: union siblings match fixed occurrence/fact footprints, fallback siblings
match in order, and negative references retain original branch results.
Unannotated unions require exact structural wiring.
\begin{mathpar}
\Rule{E-Constraint}{\vars(C)\subseteq\vars(R)\text{ for a complete earlier binder}\\ \text{flat optimizing elaboration; no let-supplied column}}
 {C\text{ may select }R\text{ before joining};\quad
  \text{otherwise select the bound prefix}}

\Rule{CT-Constraint}{\vars(C)=\bar x\quad\mathsf{columns}(\bar x)=\bar t}
 {\mathsf{facts}\longmapsto\mathsf{facts}\cup\{\mathsf{predicate}_{shape(C)}(\bar t)\}}

\Rule{RA-Constraint}{\vars(\theta)\subseteq V\quad
 \theta\text{ central and Eq/Ord-admissible}}
 {\sigma_\theta:\operatorname{Set}(V)\longrightarrow\operatorname{Set}(V)}

\Rule{CK-Elaborate}{\mathsf{WF}_{VO}(s)\quad\mathsf{WF}_{RA}(t)\\
 t.declarations=s.cq.declarations\quad t.head=s.cq.head\\
 \mathsf{FixedFacts}(s,t)\quad\mathsf{RustTree}(s,t)\\
 \forall i.\ \operatorname{Before}_s(i)\subseteq\operatorname{Before}_t(i)\\
 \mathsf{OutputBindings}(s)=\mathsf{OutputBindings}(t)}
 {\mathsf{check}_{elaborate}(s,t,())\Downarrow\ok}
\end{mathpar}
These are sufficient structural refinement obligations under the semantic
premises, not an equivalence decision or a proof of arbitrary Rust semantics.
There is no search, homomorphism construction or implication solver.
Correspondence does not repair chain order by sorting.

\paragraph{Index demands.}
The independent checker traverses RA definitions using IX-Join, IX-Const, IX-Range, IX-Other,
IX-New and IX-Reuse. It requires the next first-use descriptor at the next target position and its
normalized name \code{iN}; extra entries are rejected.
\[
\Rule{CK-Index}{\mathsf{WF}_{RA}(s)\quad\mathsf{WF}_{IX}(t)\quad t.ra=s\\
 t.indexes=[i_n=d_n]_{n=0}^{k-1}\text{ from first-use derivation}}
 {\mathsf{check}_{index}(s,t,())\Downarrow\ok}.
\]

\paragraph{Physical lowering.}
Erase probe/merge to natural join, semi to semijoin, point/range to select, and
remove modes/canon flags, preserving definitions, names, declarations and
output. The checker accepts admissible recorded decisions independently of the
default constructor. Point and range reads carry their exact const/range
supplies; central disjunctions remain filters. Let/pipe/exists choices are
checked by local use and heading premises. Anti erases to its explicit-key
logical anti and union keeps its branch origins and operands. Definition canon
boundaries, inferred keys or \code{return canon} must establish canonical
answers; host and built types establish no plan key.
\begin{mathpar}
\Rule{CK-Demands}{\mathsf{WF}_{Physical}(t)\quad
 \forall(r,C,supply)\in\mathsf{keyedReads}(t)\text{ (excluding merge-left)}.\\
 (r,\operatorname{ordinals}_{H(r)}(C),supply)\in\operatorname{descs}(s)}
 {\mathsf{demands}(s,t)\Downarrow\ok}

\Rule{PH-Range}{\theta\text{ constant order bounds on one column }c}
 {\mathsf{reads}(range\ select\ r\ if\ (\theta))=(r,\{c\},range)}

\Rule{CK-Physical}{\mathsf{WF}_{IX}(s)\quad\mathsf{WF}_{Physical}(t)\\
 \mathsf{erase}(t)=s.ra\quad\mathsf{demands}(s,t)\quad
 \mathsf{unique}(t.returned)\text{ inferred or explicit return canon}}
 {\mathsf{check}_{physical}(s,t)\Downarrow\ok}
\end{mathpar}
Demand membership includes exact stream probe/semi/anti reads, const point reads,
and range reads. Merge has a separate ordering premise. No missing IX demand
is synthesized. The checker does not choose join strategy, canonical-prefix
placement, intersection, or eagerness by rerunning lowering.

\paragraph{Structure selection.}
The default selection constructor copies Physical unchanged and records concrete types for built
key levels, each host reader's fallback, retained collections, and canonical
boundary/computed-head dedup state. The target language independently checks
read coverage, key levels, entry identity, known representation compatibility,
and canonicality. Arbitrary aliases or user implementations remain subject to
the generated Rust trait obligations; a valid target need not use default types.
Compatible membership prefixes may share a curried structure. A recorded
smallest intersection requires two candidate levels on every retained side;
reusable hosts instead record candidate and membership fallbacks. Const/range
supplies never justify index construction. Native ordered host reads and
same-key merge pipes discharge their written ordering obligations. Entry order is free; checking compares no constructor token stream. The
selection correspondence accepts any locally valid target whose embedded
Physical satisfies \code{physical::check_rewrite} against the source, so it
can also change admissible physical choices.
\[
\Rule{CK-Structures}{\mathsf{WF}_{Physical}(s)\quad\mathsf{WF}_{Structures}(t)\\
 \mathsf{checkRewrite}_{Physical}(s,t.physical)\Downarrow\ok}
 {\mathsf{check}_{structures}(s,t)\Downarrow\ok}.
\]

\subsection{Realization evidence}
\code{cq-physical-realization} checks its Structures input and prints a
\code{syn::Expr}; it has no source/target \code{check} and no independent Rust
semantic checker. Its evidence consists of independent answer oracles, rustc
trait guards at generated uses, and authored expansion fixtures
(Section~\ref{sec:evaluation}). The optional binary intersection path chooses
between candidate sets using runtime counts without reordering the RA/VO plan.
Rows borrow host/built fields, computed columns remain owned, and host lookups
construct owned lookup keys; the recorded dedup type determines row retention. A block
with bound storage collects answers before returning \code{Vec::into_iter};
a block without bound statements may return the lazy iterator. This evidence
supports the exercised cases, not an arbitrary Rust semantic equality theorem.

\subsection{The four interfaces of the closed core}
Extensions compose checked query invocations in ordinary Rust contexts;
each invocation passes through the same closed core. Surface extensions introduce no operators.
The following rules are implementation obligations, not proofs of extension
semantics or Rust expression semantic equality.
\begin{mathpar}
\Rule{IF-Surface}{\operatorname{parse}_{CQ}(T)=s\quad
 \text{no unexpanded extension clause remains in }s}
 {\mathsf{surface}(T)=s\text{ is admitted to CQ checking}}

\Rule{IF-Input}{\text{unique relation and column names}\quad\text{positive arity}}
 {\mathsf{layout}(d)=\operatorname{columns}(d)\text{ in declaration order}}

\Rule{IF-Output}{A=\mathsf{collect\_from}(I)\quad
 \alpha(A)=\{r\mid r\in I\}}
 {A:\code{Canonical}\quad A\text{ may supply an input host}}

\Rule{IF-Environment}{\mathsf{visibleLogicalUses}(e)\subseteq\Gamma_Q\quad
 \Gamma_H,\Gamma_Q\vdash_{Rust}e}
 {e\text{ is embedded Rust in the closing environment}}

\Rule{IF-Parameter}{p\in\vars_H(\theta)\quad b_p\text{ fresh}}
 {p\longmapsto\code{let b_p = \&p;}\text{ once outside all query loops}}

\Rule{IF-Range}{C=\{d_1\}\quad H:\code{OrderedRead}\text{ at the matching key type}}
 {\mathsf{window}(H,C)\text{ is an ordered read; otherwise scan/filter, no build}}

\Rule{IF-Leaf}{d=\code{relation R(...)=}(e)\quad b\text{ fresh}\quad
 \Gamma_H\vdash_{Rust}e}
 {d\longmapsto\code{let b = \&(e);}\quad R\mapsto b}

\Rule{IF-Dispatch}{C=\{d_1,\ldots,d_k\}\quad
 H:\mathsf{ReadAt}(\tau_1,\ldots,\tau_k)}
 {\mathsf{read}(H,C)\longmapsto\mathsf{nativeRead}(H,C)}

\Rule{IF-Fallback}{C\text{ is nonleading}\ \lor\ H\not:\mathsf{ReadAt}(C)}
 {\mathsf{read}(H,C)\longmapsto\mathsf{lazyOnce}(\mathsf{build}(\alpha(H),C))}
\end{mathpar}
Stable autoref specialization chooses a native keyed read only at the exact
leading-prefix key type. Stream reads otherwise allocate one lazy fallback per key; const/range
reads scan and filter without building. Range uses the leading column when
its host type orders that key. Copied bound keys permit native BTreeMap
lower-bound range seeking; other ordered cursors step to the bound using
borrowed values without invoking Clone. Inverted bounds yield no rows. Merge is direct-only on a leading host key and requires
OrderedRead, diagnosed by rustc. Canonicality never follows from a host type.
Grouped answers concatenate the leading key tuple with the remaining columns;
ungrouped answers supply full rows. Both may retain references. Leaf bindings
occur once per query evaluation in declaration order. Pass checks preserve
complete declarations and leaf syntax without evaluating Rust expressions.
